\DocumentMetadata{
  pdfversion=2.0,
  pdfstandard=ua-2,
  lang=en-US,
  tagging=on,
  tagging-setup={math/setup=mathml-SE,tabsorder=structure}
}
\documentclass[11pt,letterpaper]{article}
\usepackage[margin=0.68in,includefoot]{geometry}
\usepackage{newcomputermodern}
\usepackage{amsmath,amscd,mathtools}
\usepackage{tikz,adjustbox}
\providecommand{\square}{\mdlgwhtsquare}
\usepackage[hidelinks]{hyperref}
\hypersetup{
  pdftitle={Analysis PhD Exam, September 2005},
  pdfauthor={University of Florida Department of Mathematics},
  pdfsubject={Graduate mathematics examination},
  pdfdisplaydoctitle=true,
  bookmarksopen=true
}
\setcounter{secnumdepth}{0}
\setlength{\parindent}{0pt}
\setlength{\parskip}{0.28em}
\setlength{\emergencystretch}{3em}
\setlength{\leftmargini}{2.15em}
\setlength{\leftmarginii}{2.65em}
\setlength{\leftmarginiii}{3em}
\pagestyle{empty}
\raggedbottom
\allowdisplaybreaks
\NewTaggingSocketPlug{block/list/label}{examproblem}{%
  \tagstructbegin{tag=Lbl}\tagstructend%
  \tagstructbegin{tag=\LBody}%
  \tagstructbegin{tag=\ProblemHeadingTag,title-o={\ProblemHeadingText}}%
  \tagmcbegin{tag=\ProblemHeadingTag,actualtext={\ProblemHeadingText}}%
  #2%
  \tagmcend%
  \tagstructend%
}
\newenvironment{examproblems}[2]{%
  \def\ProblemHeadingTag{#2}%
  \AssignTaggingSocketPlug{block/list/label}{examproblem}%
  \begin{enumerate}%
  \setcounter{enumi}{#1}%
  \renewcommand{\labelenumi}{\textbf{\theenumi.}}%
  \setlength{\topsep}{0.35em}%
  \setlength{\partopsep}{0pt}%
  \setlength{\itemsep}{0.48em}%
  \setlength{\parsep}{0.18em}%
}{\end{enumerate}}
\newenvironment{examsubparts}{%
  \AssignTaggingSocketPlug{block/list/label}{default}%
  \begin{enumerate}%
  \setlength{\topsep}{0.18em}%
  \setlength{\partopsep}{0pt}%
  \setlength{\itemsep}{0.16em}%
  \setlength{\parsep}{0.1em}%
}{\end{enumerate}}
\newenvironment{examinstructions}{%
  \par\begingroup\itshape
}{%
  \par\endgroup\vspace{0.18em}
}
\makeatletter
\renewcommand\subsection{\@startsection{subsection}{2}{0pt}%
  {-1.25ex plus -.25ex minus -.1ex}%
  {0.55ex plus .1ex}%
  {\normalfont\large\bfseries}}
\renewcommand{\@maketitle}{%
  \newpage\null\vskip -1em%
  {\footnotesize\bfseries
    \href{https://www.ufl.edu/}{University of Florida}, \href{https://math.ufl.edu/}{Department of Mathematics}\par}%
  \vskip -0.5em%
  \tagstructbegin{tag=H1,title={Analysis PhD Exam, September 2005}}%
  \tagpdfparaOff%
  \tagmcbegin{tag=H1}%
    {\LARGE\bfseries\href{https://gma.math.ufl.edu/exams/analysis-phd/}{Analysis PhD Exam}, September 2005\par}%
  \tagmcend%
  \tagpdfparaOn%
  \tagstructend%
  \vskip 0.25em%
  \tagmcbegin{artifact}%
    \hrule height 1pt%
  \tagmcend%
  \vskip 1em%
}
\makeatother
\title{Analysis PhD Exam, September 2005}
\author{}
\date{}
\begin{document}
\maketitle
\thispagestyle{empty}
\begin{examinstructions}
Give complete proofs and computations. Be particularly careful to indicate why the most crucial steps in your proof are true. Partial credit will be given where justified.
\end{examinstructions}
\begin{examproblems}{0}{H2}
\def\ProblemHeadingText{Problem 1.}
\item
Let \(\Omega\subset\mathbb R^n\) be open. Prove that \(A\subset\mathcal D(\Omega)\) is bounded if and only if
\[
\sup\{|\Lambda(\varphi)|\mid \varphi\in A\}<\infty,
\]
 for every \(\Lambda\in\mathcal D'(\Omega)\).
\def\ProblemHeadingText{Problem 2.}
\item
Let~\(\mu\) be Lebesgue measure on \(\mathbb R\) and~\(f\) be~\(\mu\)-integrable on \([0,1]\). Suppose \(\int_0^r f\,d\mu=0\), for every rational \(r\in[0,1]\). Give a complete characterization of~\(f\).
\def\ProblemHeadingText{Problem 3.}
\item
Let \(\mathcal H\) be a Hilbert space and suppose~\(A\) is a bounded operator on \(\mathcal H\) such that the range of~\(A\) is finite dimensional. Let \(\{y_1,\ldots,y_n\}\) be an orthonormal basis for the range of~\(A\). Prove that there exist \(x_1,\ldots,x_n\in\mathcal H\) such that
\[
Ax=\sum_{j=1}^n[x,x_j]y_j,
\]
 for all \(x\in\mathcal H\).
\def\ProblemHeadingText{Problem 4.}
\item
Let~\(A\) be a bounded operator on a Hilbert space \(\mathcal H\) such that \(A(B_1)=\{Ax\mid x\in B_1\}\) is compact, where \(B_1\) is the closed unit ball of \(\mathcal H\). Prove that, as a map from \(B_1\) with the weak vector topology to \(\mathcal H\) with the norm topology, \(A\) is continuous. (Hint: show that the weak and norm topologies on \(A(B_1)\) must coincide under the given hypothesis.)
\def\ProblemHeadingText{Problem 5.}
\item
Let~\(\mu\) be a positive Radon measure on \(\mathbb R^N\) and~\(f\) be~\(\mu\)-integrable. Show that for each \(\epsilon>0\) there exist an upper semi-continuous function \(u_\epsilon\) and a lower semi-continuous function \(v_\epsilon\) such that \(u_\epsilon\leq f\leq v_\epsilon\) and \(\int(v_\epsilon-u_\epsilon)\,d\mu<\epsilon\).
\def\ProblemHeadingText{Problem 6.}
\item
Let~\(\mu\) be a finite positive Radon measure on \(\mathbb R^N\) and \(\{f_n\}_{n\in\mathbb N}\subset L^2(\mu)\) with \(\|f_n\|_2\leq1\). Suppose \(\{f_n\}_{n\in\mathbb N}\) converges~\(\mu\)-a.e. to~\(f\). Show that \(\int|f_n-f|\,d\mu\to0\), as \(n\to\infty\).
\def\ProblemHeadingText{Problem 7.}
\item
Let~\(E\) be a topological vector space. Prove that if~\(A\) and~\(B\) are compact subsets of~\(E\), then so is \(A+B\).
\def\ProblemHeadingText{Problem 8.}
\item
Consider \(X\in\mathcal D(\mathbb R^2)^*\) defined by
\[
\mathcal D(\mathbb R^2)\ni\phi\mapsto X(\phi)=\int_{\mathbb R}\phi(x,-x)\,dx.
\]
 Note. So~\(X\) is a solution of the PDE \((\partial_1-\partial_2)X=0\) which is not a classical solution.
\begin{examsubparts}
\item[\textnormal{(a)}]
Show that~\(X\) is a distribution of order zero.
\item[\textnormal{(b)}]
What is the support of~\(X\)? (Prove it.) Use this to prove that~\(X\) is not a regular distribution with continuous coefficient function.
\item[\textnormal{(c)}]
Show that \((\partial_1-\partial_2)X=0\).
\end{examsubparts}
\end{examproblems}
\end{document}
